\documentclass[11pt]{article}

\usepackage[margin=1in]{geometry}
\usepackage{amsmath,amssymb,amsthm,mathtools,mathrsfs}
\usepackage{booktabs,microtype,xurl,xcolor}
\usepackage[colorlinks=true,linkcolor=blue!55!black,
            citecolor=blue!55!black,urlcolor=blue!55!black]{hyperref}

\newtheorem{theorem}{Theorem}[section]
\newtheorem{proposition}[theorem]{Proposition}
\newtheorem{corollary}[theorem]{Corollary}
\newtheorem{lemma}[theorem]{Lemma}
\theoremstyle{remark}
\newtheorem{remark}[theorem]{Remark}

\newcommand{\C}{\mathbb C}
\newcommand{\End}{\operatorname{End}}
\newcommand{\Smith}{\operatorname{Smith}}

\title{A Non-Split Regge--Wheeler Self-Extension and a\\
Defective Schwarzschild Resonance in Pure Weyl Gravity}
\author{Authors withheld for review}
\date{26 July 2026}

\begin{document}
\maketitle

\begin{abstract}
We determine the axial Schwarzschild spin-two extension and one of its
quasinormal resonances in four-dimensional pure Weyl gravity.  The reduced
six-state Bach system has an exact filtration with two spin-two
Regge--Wheeler factors and one spin-one Maxwell factor.  The repeated
spin-two block is not a rational direct sum: it is a non-split
self-extension and the first parameter jet of a lower-dimensional family.
After Liouville reduction, its projective class is represented by
\[
 [\mathcal I_{\rm Bach}]
 =\frac{i\omega}{2}\left[1-\frac2r\right].
\]
The class in brackets is exactly the transverse-traceless critical
Einstein--Weyl mass direction.  The required rational gauge has the
linearly growing asymptotics of the differentiated massive Jost phase;
the long-range Coulomb contribution cancels the apparent logarithmic
phase derivative.  Thus the local Bach jet and the differentiated
physical Jost germs agree, up to endpoint normalization.

At a validated simple Schwarzschild quasinormal frequency, the connection
shear is nonzero and the spin-one factor is a local unit.  The complete
connection therefore has Smith type \((0,0,2)\).  Its geometric root is
Einstein, while the generalized root has a nonzero gauge-invariant
Ricci-carrier component.  The same invariant is the nonzero velocity of
the massive spin-two quasinormal branch.  We prove that the compact-source,
compact-observation outgoing radial Green operator on the Schwarzschild
exterior has a nonzero rank-one second-order pole.

Exact axial reconstruction shows that the Einstein range of the leading
pole has nonzero Bondi shear, and an explicit conserved, traceless,
compactly supported odd source has nonzero adjoint overlap.  Hence both
source and observation channels can be nonvanishing.  The isolated
resonance contour contains the expected polynomially weighted Einstein
profile.  We do not claim a global retarded quasinormal expansion: the
time-independent generalized component has weakened \(O(1)\) strain, and
a closed weighted Fredholm realization with causal contour control remains
open.
\end{abstract}

\section{Introduction}

Fourth-order linear equations are often described by listing their scalar
factors.  That description is incomplete when a factor is repeated.  The
layers may form a direct sum, or they may form a non-semisimple extension.
Only the second possibility carries an intrinsic generalized state and can
turn a simple resonance of the repeated factor into a defective resonance
of the full system.

This distinction is realized by axial Schwarzschild perturbations of pure
Weyl gravity.  In the factor-adapted radial variables, the Bach equation
contains two spin-two Regge--Wheeler layers and one spin-one Maxwell layer,
but the two spin-two layers cannot be separated by a rational local field
redefinition.  The extension is also the critical limit of the massless
and massive spin-two branches of Einstein--Weyl gravity.  At a
Schwarzschild quasinormal frequency, this confluence produces a
length-two root chain rather than two independent Einstein roots.

The paper has five results.
\begin{enumerate}
\item The axial spin-two Bach block is a non-split rational
  Regge--Wheeler self-extension and an exact first parameter jet.
\item Its projective cocycle has the elementary mass-direction normal form
  \[
    [\mathcal I_{\rm Bach}]
    =\frac{i\omega}{2}\left[1-\frac2r\right],
  \]
  and the equivalence extends to differentiated horizon and infinity Jost
  germs.
\item A validated QNM calculation selects the defective Smith type
  \((0,0,2)\).  The generalized root necessarily crosses the Einstein
  filtration and is the normalization-independent tangent class of the
  massive QNM.
\item The compact-source, compact-observation exterior outgoing radial
  Green operator has a genuine rank-one pole of order two.
\item Exact metric reconstruction and an explicit conserved, traceless odd
  source show that the range and source covectors of this pole can both be
  nonzero.
\end{enumerate}

These statements combine exact algebra with validated numerics.  The
module reduction, cocycle identities, mass comparison, Green reduction,
metric reconstruction, and source construction are exact.  The existence
and uniqueness of the enclosed QNM, nonvanishing of the projective
selector, and spin-one local-unit condition are interval-validated.

\subsection*{Claim boundary}

\begin{center}
\footnotesize
\begin{tabular}{p{0.34\linewidth}p{0.36\linewidth}p{0.20\linewidth}}
\toprule
Result & Input & Status\\
\midrule
Filtration, partial jet, and nonsplitting
& Rational identities and exhaustive coboundary obstruction
& \texttt{LOCAL-ALGEBRAIC}\\
Mass-direction and Jost-germ comparison
& Exact cocycle primitive and moving massive phases
& \texttt{LOCAL-ALGEBRAIC}\\
Defective QNM and non-Einstein root
& Validated Evans winding, selector, spin-one unit
& \texttt{REDUCED-MODE}\\
Exterior cut-off Green double pole
& Analytic Jost maps and exact transparent cuts
& \texttt{REDUCED-MODE}\\
Bondi-shear and compact-source nonannihilation
& Exact reconstruction and conserved-source identities
& \texttt{REDUCED-MODE}\\
Global retarded ringdown expansion
& Weighted Fredholm domains and contour deformation
& Not established\\
\bottomrule
\end{tabular}
\end{center}

The last row is essential.  The paper proves a meromorphic exterior radial
response after compact source and observation localization, and it proves
the polynomial contribution of an isolated local resonance contour.  It
does not prove a complete causal quasinormal expansion, late-time
asymptotics, detector sensitivity, or excitation by a specified
astrophysical matter model.

\section{The filtered axial Bach module}
\label{sec:filtration}

\subsection{Differential-module conventions}

We work over
\[
 K=\C(r,\omega),\qquad
 D=f\partial_r,\qquad
 f=1-\frac2r,
\]
with \(\omega\ne0\) unless stated otherwise.  The spin-two and spin-one
scalar operators for axial \(\ell=2\) are
\[
 L_2=D^2+U_2,\qquad
 U_2=\omega^2-f\left(\frac6{r^2}-\frac6{r^3}\right),
\]
\[
 L_1=D^2+U_1,\qquad
 U_1=\omega^2-f\frac6{r^2}.
\]
We use \(D(q)\) and \(D(U)\), rather than primes, for the radial
derivation.  A prime on a connection coefficient below denotes
\(\partial_\omega\).

The Einstein submodule is the kernel of the gauge-invariant
traceless-Ricci, equivalently Schouten, carrier.  Thus ``non-Einstein''
means nonzero projection to this carrier quotient, not failure of a
coordinate gauge choice.

\subsection{Filtration and partial jet}

In the certified factor gauge, the six-state system is
\begin{equation}\label{eq:six-state}
\begin{aligned}
 DX&=AX+EY+CZ,\\
 DY&=AY+DZ,\\
 DZ&=A_1Z,
\end{aligned}
\end{equation}
where \(X,Y,Z\in K^2\), while \(A\) and \(A_1\) are the companion
matrices of \(L_2\) and \(L_1\).  The rational reconstruction into the
axial metric and curvature variables has full row rank on the declared
domain.

\begin{theorem}[Filtered axial Bach module]\label{thm:filtration}
The reduced axial \(\ell=2\) Bach module has a filtration with successive
quotients
\[
 M_2,\qquad M_2,\qquad M_1.
\]
The middle four-state object is an extension
\[
 0\longrightarrow M_2\longrightarrow E_2
 \longrightarrow M_2\longrightarrow0,
\]
and the final quotient is the spin-one Maxwell module \(M_1\).
\end{theorem}

This theorem is an exact rational statement.  It does not rely on an
endpoint normalization or on the QNM calculation.

Introduce the four-state family
\[
 \mathcal B(\tau)=
 \begin{pmatrix}
 A+\tau E&D+\tau C\\
 0&A_1
 \end{pmatrix},
\qquad
 \Phi_4(\tau)=
 \begin{pmatrix}P(\tau)&Q(\tau)\\0&R\end{pmatrix}.
\]

\begin{theorem}[Exact partial jet]\label{thm:partial-jet}
The system \eqref{eq:six-state} is the first spin-two-row jet of
\(\mathcal B(\tau)\) at \(\tau=0\).  Its fundamental matrix is
\[
 \Phi_6=
 \begin{pmatrix}
 P&\dot P&\dot Q\\
 0&P&Q\\
 0&0&R
 \end{pmatrix}_{\tau=0},
\qquad
 \det\Phi_6=(\det P)^2\det R.
\]
Every factor-adapted endpoint connection consequently has the form
\begin{equation}\label{eq:triangular-connection}
 T(\omega)=
 \begin{pmatrix}
 a&b&c\\
 0&a&d\\
 0&0&g
 \end{pmatrix},
\qquad
 \det T=a^2g.
\end{equation}
\end{theorem}

\begin{proof}
Differentiate
\(D\Phi_4(\tau)=\mathcal B(\tau)\Phi_4(\tau)\) at \(\tau=0\).
The differentiated first row, the undifferentiated spin-two row, and the
spin-one row reproduce \eqref{eq:six-state}.  The determinant and
connection identities follow from block triangularity.
\end{proof}

\subsection{Projective extension class}

Consider a scalar self-extension
\[
 Lx=(s_1D+s_0)y,\qquad Ly=0,\qquad L=D^2+U.
\]
Define
\[
 \mathcal I=s_0-\frac12D(s_1),\qquad
 \mathcal K_U=D^3+4UD+2D(U).
\]

\begin{lemma}[Triangular gauge]\label{lem:triangular-gauge}
For \(q\in K\), set
\[
 Q(q)=qD-\frac12D(q).
\]
On \(\ker L\),
\[
 [L,Q(q)]=-\frac12\mathcal K_U(q).
\]
Thus the lift change \(x\mapsto x+Q(q)y\) changes the scalar density by
\[
 \mathcal I\longmapsto
 \mathcal I-\frac12\mathcal K_U(q).
\]
\end{lemma}

\begin{proof}
Expand the commutator by the Leibniz rule and use \(D^2y=-Uy\).
The coefficient of \(Dy\) cancels; the remaining coefficient is
\(-[D^3q+4UDq+2D(U)q]/2\).
\end{proof}

The projective extension class is therefore
\[
 [\mathcal I]\in K/\mathcal K_UK.
\]
The operator \(\mathcal K_U\) is the symmetric-square operator of \(L\).

\begin{theorem}[Mass-direction normal form]\label{thm:mass-normal}
For the axial \(\ell=2\) Bach extension,
\[
 \mathcal I_{\rm Bach}
 =
 \frac{i(r-2)(2r\omega^2+3\omega^2+12)}
 {5r^4\omega}.
\]
Let
\[
 q(r,\omega)=
 -\frac{i}{120\omega}
 \left(15r+13+\frac{12}{r}+\frac9{r^2}\right).
\]
Then, for \(\omega\ne0\),
\begin{equation}\label{eq:mass-normal-form}
 \mathcal I_{\rm Bach}
 =\mathcal K_{U_2}q+\frac{i\omega}{2}f,
\end{equation}
and hence
\[
 \boxed{\;
 [\mathcal I_{\rm Bach}]
 =\frac{i\omega}{2}[f].
 \;}
\]
\end{theorem}

\begin{proof}
Substitute \(D=f\partial_r\), \(U_2\), and the displayed \(q\) into
\(\mathcal K_{U_2}q\).  After clearing \(120r^4\omega\), the assertion is
a polynomial identity in \(r\) and \(\omega\).  Exact coefficient
comparison gives \eqref{eq:mass-normal-form}.
\end{proof}

\begin{theorem}[Non-split physical-axis extension]\label{thm:nonsplit}
For every real \(\omega>0\), the extension \(E_2\) is non-split over
\(\C(r)[D]\).
\end{theorem}

\begin{proof}[Certificate-supported exact proof]
The rational splitting equation has exhaustive pole and degree bounds and
reduces to a finite coefficient system.  Fixed minors of its coefficient
and augmented matrices are
\[
 \det M_{[0,1,2]}=3456\,\omega^{10},
 \qquad
 \det[M\mid{\rm rhs}]_{[0,1,2,5]}
 =-645120\,i\,\omega^9.
\]
They are simultaneously nonzero for every real \(\omega>0\), so the
inhomogeneous coboundary equation is inconsistent.
\end{proof}

Theorem~\ref{thm:nonsplit} is a statement about the differential module;
Theorem~\ref{thm:mass-normal} identifies a simple representative of its
class.  Neither theorem asserts that every QNM of the repeated scalar
factor must be defective.  Defectiveness is a resonant evaluation of the
nonzero extension class.

\section{Critical-mass interpretation}
\label{sec:mass}

\subsection{The transverse-traceless mass jet}

On the transverse-traceless spin-two sector, introduce the
Einstein--Weyl family
\[
 (\mathcal E+m\mathcal A)\phi=0.
\]
Its axial scalar equation is
\begin{equation}\label{eq:massive-rw}
 L_m(\omega)y=
 \left[D^2+\omega^2-V_2-mf\right]y=0.
\end{equation}
For a differentiable solution family,
\[
 L_2\,\partial_my=f\,y
 \qquad (m=0).
\]

\begin{theorem}[Exact critical-mass jet]\label{thm:mass-jet}
The projective class of the transverse-traceless mass jet is
\[
 [\mathcal I_{\rm mass}]=[f].
\]
Consequently,
\[
 \boxed{\;
 [\mathcal I_{\rm Bach}]
 =\frac{i\omega}{2}[\mathcal I_{\rm mass}],
 \qquad
 m=\frac{i\omega}{2}\tau.
 \;}
\]
The parameter identity is understood to first order and modulo rational
triangular gauge.
\end{theorem}

\begin{proof}
Differentiating \eqref{eq:massive-rw} gives the extension density \(f\).
The result follows from Theorem~\ref{thm:mass-normal}.
\end{proof}

\subsection{Endpoint-compatible gauge}

Because Lemma~\ref{lem:triangular-gauge} contains a factor \(1/2\), the
field redefinition that removes \(\mathcal K_{U_2}q\) is
\[
 Q_q=2qD-D(q).
\]
On \(\ker L_2\),
\[
 [L_2,Q_q]=-\mathcal K_{U_2}(q).
\]
Thus, if \(L_2x=\mathcal I_{\rm Bach}y\), then
\[
 \widehat x=x+Q_qy
 \quad\Longrightarrow\quad
 L_2\widehat x=\frac{i\omega}{2}fy.
\]

The primitive has the forced infinity behavior
\[
 q(r,\omega)=-\frac{i}{8\omega}r+O(1).
\]
This growth is not an endpoint obstruction.  It is the derivative of the
moving massive phase.

\begin{theorem}[Differentiated Jost compatibility]
\label{thm:jost-compatibility}
Let \(y_\sigma(m,r)\), \(\sigma=\pm1\), be a massive Jost germ at infinity
with
\[
 k=(\omega^2-m)^{1/2},\qquad
 y_\sigma(m,r)=
 e^{\sigma ikr}r^{\rho_\sigma(m)}
 \left(1+O(r^{-1})\right),
\]
\[
 \rho_\sigma(m)
 =\sigma i\left(2k+\frac{m}{k}\right).
\]
Then
\[
 \rho_\sigma'(0)=0,\qquad
 \frac{\partial_my_\sigma(0,r)}{y_\sigma(0,r)}
 =-\frac{\sigma i}{2\omega}r+O(1).
\]
With \(c(\omega)=i\omega/2\),
\[
 \frac{Q_qy_\sigma}{y_\sigma}
 =\frac{\sigma r}{4}+O(1)
 =c(\omega)
 \frac{\partial_my_\sigma(0,r)}{y_\sigma(0,r)}+O(1).
\]
At the horizon, the mass perturbation is multiplied by \(f=O(r-2)\),
so the indicial exponent is unchanged and \(Q_q\) preserves the analytic
ingoing germ.  At each endpoint \(e\),
\[
 \widehat x_e-c(\omega)\partial_my_e(0)
 =\lambda_e(\omega)y_e
\]
for an analytic scalar normalization \(\lambda_e\).
\end{theorem}

\begin{proof}
The infinity equation is
\[
 D^2y+\left(k^2+\frac{2m}{r}+O(r^{-2})\right)y=0.
\]
The coefficient of \(r^{-1}\) gives
\[
 2\sigma ik\rho_\sigma+4k^2+2m=0.
\]
Since \(k'(0)=-1/(2\omega)\),
\[
 \left.\frac{d}{dm}\left(2k+\frac{m}{k}\right)\right|_{m=0}
 =-\frac1\omega+\frac1\omega=0.
\]
The Coulomb-power derivative therefore cancels, while differentiation of
the plane-wave phase gives the linear term.  Finally,
\(Dy_\sigma/y_\sigma=\sigma i\omega+O(r^{-1})\) and
\(q=-ir/(8\omega)+O(1)\) give the stated asymptotics of \(Q_qy_\sigma\).
The difference between \(\widehat x\) and
\(c\,\partial_my\) is homogeneous and belongs to the same one-dimensional
endpoint Jost class.  At the horizon the claim follows from regular
Frobenius dependence and the unchanged indicial exponent.
\end{proof}

\begin{corollary}[Evans derivative and mass velocity]
\label{cor:mass-velocity}
Let \(a(\omega,m)\) be the massive spin-two Evans coefficient in analytic
Jost normalizations, and let \(b_{\rm B}(\omega)\) be the Bach connection
shear.  There is an analytic \(h(\omega)\) such that
\[
 b_{\rm B}(\omega)
 =\frac{i\omega}{2}\partial_ma(\omega,0)
 +h(\omega)a(\omega,0).
\]
At a simple QNM \(\omega_n\),
\[
 b_{\rm B}(\omega_n)
 =\frac{i\omega_n}{2}\partial_ma(\omega_n,0).
\]
Writing
\[
 \kappa_n=\frac{b_{\rm B}(\omega_n)}
 {a_\omega(\omega_n,0)},
\]
the massive QNM branch satisfies
\begin{equation}\label{eq:mass-velocity}
 \boxed{\;
 \omega_n'(0)=\frac{2i}{\omega_n}\kappa_n.
 \;}
\end{equation}
\end{corollary}

\begin{proof}
The endpoint corrections in
Theorem~\ref{thm:jost-compatibility} are scalar changes of differentiated
Jost normalizations, hence contribute an analytic multiple of \(a\).
At a zero of \(a\), this term vanishes.  Differentiate
\(a(\omega_n(m),m)=0\) and solve for \(\omega_n'(0)\).
\end{proof}

This is the asymptotically flat QNM analogue of the mass-derivative
construction of logarithmic partners in critical gravity
\cite{lupope2011,bergshoeff2011}.  The terminology must be used with care:
the Schwarzschild endpoint tangent is polynomial, not an independent
radial \(\log r\) mode.

\section{The defective Schwarzschild resonance}
\label{sec:qnm}

\subsection{Smith dichotomy}

At a simple spin-two zero, analytic elimination of the spin-one unit
reduces \eqref{eq:triangular-connection} to
\[
 T_2(\omega)=
 \begin{pmatrix}
 a(\omega)&b(\omega)\\
 0&a(\omega)
 \end{pmatrix},
\qquad
 a(\omega_n)=0,\quad a'(\omega_n)\ne0.
\]

\begin{proposition}[Local Smith dichotomy]\label{prop:smith}
If \(g(\omega_n)\ne0\), then
\[
 \Smith(T)=
 \begin{cases}
 (0,0,2),&b(\omega_n)\ne0,\\
 (0,1,1),&b(\omega_n)=0.
 \end{cases}
\]
The first case has one length-two root chain; the second has two
independent length-one roots.
\end{proposition}

\begin{proof}
With \(z=\omega-\omega_n\),
\(a=a_1z+O(z^2)\), \(a_1\ne0\).  The determinant valuation of the
spin-two block is two.  If \(b(\omega_n)\ne0\), its first Fitting ideal is
the unit ideal and the exponents are \((0,2)\).  If \(b(\omega_n)=0\),
all entries vanish to first order and the exponents are \((1,1)\).
The spin-one unit contributes a zero exponent.
\end{proof}

\subsection{Validated specialization}

The positively oriented validated contour encloses
\[
 \Omega_n=
 \left\{\omega:
 \left|\omega-
 \bigl(-0.3736716844180418+0.0889623156889357\,i\bigr)
 \right|\le10^{-7}\right\}.
\]
Projective interval panels exclude zeros on the contour and give winding
number \(+1\).  The disk therefore contains exactly one simple spin-two
zero.  On the localized disk,
\[
 \delta_\tau(\Omega_n)
 \subset
 [0.0010932\pm7.60\times10^{-8}]
 +i[-0.0002195\pm8.89\times10^{-8}],
\]
which excludes zero.  An independent local calculation bounds the
spin-one mismatch modulus below by \(1/2500\).

\begin{theorem}[Certified defective resonance]\label{thm:certified-ep2}
At the unique QNM in \(\Omega_n\), the complete axial connection has
Smith valuations
\[
 \boxed{\;(0,0,2).\;}
\]
Its spin-two divisor has algebraic multiplicity two, geometric
multiplicity one, and one length-two root chain.
\end{theorem}

\begin{proof}
The winding certificate supplies the simple spin-two zero.  The selector
enclosure gives \(b(\omega_n)\ne0\), and the spin-one bound gives
\(g(\omega_n)\ne0\).  Proposition~\ref{prop:smith} applies.
\end{proof}

\subsection{One resonant invariant}

Encode the endpoint conditions in an analytic finite-interval pencil
\[
 L(\omega):X\longrightarrow Y.
\]
Choose right and left states at the simple scalar resonance,
\[
 L_nu_n=0,\qquad L_n^*\widetilde u_n=0,
\]
and define, with all augmented endpoint contributions included,
\[
 \alpha_n=
 \langle\widetilde u_n,L_n'u_n\rangle\ne0.
\]
For the repeated block
\[
 \mathbb L(\omega)=
 \begin{pmatrix}L(\omega)&K(\omega)\\0&L(\omega)\end{pmatrix},
\]
set
\[
 \beta_n=\langle\widetilde u_n,K_nu_n\rangle.
\]

\begin{theorem}[Resonant evaluation]\label{thm:resonant-evaluation}
In endpoint frames compatible with the partial jet,
\begin{equation}\label{eq:unified-invariant}
 \boxed{\;
 \kappa_n
 =\frac{b(\omega_n)}{a'(\omega_n)}
 =\frac{\beta_n}{\alpha_n}
 =-\left.\frac{d\omega_n}{d\tau}\right|_0
 =-\frac{i\omega_n}{2}
 \left.\frac{d\omega_n}{dm}\right|_0.
 \;}
\end{equation}
The ratio is invariant under a common analytic nonzero Evans unit and
under triangular changes of splitting.
\end{theorem}

\begin{proof}
Differentiate
\[
 [L(\omega)+\tau K(\omega)]u(\tau)=0
\]
along the resonance branch.  Pairing with \(\widetilde u_n\) gives
\[
 \omega_n'(0)\alpha_n+\beta_n=0.
\]
Differentiating the connection divisor gives
\[
 a'(\omega_n)\omega_n'(0)+b(\omega_n)=0,
\]
because the partial jet identifies \(b=\partial_\tau a\).  The mass
identity follows from \eqref{eq:mass-velocity}.  A triangular change
\(K\mapsto K+LQ-QL\) has zero resonant pairing, and an Evans unit
multiplies numerator and denominator by the same nonzero value.
\end{proof}

\begin{proposition}[Non-Einstein generalized root]
\label{prop:non-einstein-root}
The geometric root belongs to the Einstein submodule.  Every length-two
root polynomial has generalized coefficient with nonzero projection to
the Ricci-carrier quotient.  In the normalized triangular frame, that
projection is
\[
 \boxed{\;
 -\frac{a'(\omega_n)}{b(\omega_n)}
 =-\frac{\alpha_n}{\beta_n}
 =-\frac1{\kappa_n}\ne0.
 \;}
\]
\end{proposition}

\begin{proof}
Write \(a=a_1z+O(z^2)\), \(b=b_0+O(z)\), with \(a_1b_0\ne0\).
For a root polynomial
\[
 c(z)=c_0+zc_1+O(z^2),\qquad
 T_2(z)c(z)=O(z^2),
\]
the order-zero equation gives \(c_0=e_X\), the Einstein vector.
Writing \(c_1=xe_X+ye_Y\), the order-\(z\) equation is
\[
 a_1+b_0y=0.
\]
Hence \(y=-a_1/b_0\ne0\).  Adding an Einstein vector changes \(x\), not
the carrier coefficient \(y\).
\end{proof}

Thus the defective resonance is not the same Einstein QNM counted twice.
Its generalized member is the canonical mass-tangent class modulo the
ordinary QNM.

\section{Exterior Green pole and critical response}
\label{sec:green}

\subsection{Finite cuts and the exterior inverse}

Put \(x=r_*\).  On the certified QNM neighborhood, let
\[
 u_H(x,\omega),\qquad u_\infty(x,\omega)
\]
be analytic spin-two Jost solutions satisfying the future-horizon and
outgoing-infinity conditions.  Their Wronskian is the scalar Evans
coefficient, up to an analytic unit.

Choose nontrivial smooth source and observation cutoffs
\(\chi_{\rm s},\chi_{\rm o}\) and finite cuts
\[
 x_-<
 \operatorname{supp}\chi_{\rm s}\cup
 \operatorname{supp}\chi_{\rm o}
 <x_+.
\]
Impose the exact transparent boundary conditions
\[
 W(u,u_H)|_{x_-}=0,\qquad
 W(u,u_\infty)|_{x_+}=0.
\]
They are analytic in \(\omega\), and uniqueness of the exterior Jost
problems identifies the finite-interval inverse between the cutoffs with
\(\chi_{\rm o}R_{\rm ext}(\omega)\chi_{\rm s}\).

\begin{theorem}[Exterior cut-off Green pole]
\label{thm:exterior-green-pole}
For any such nontrivial compact source and observation cutoffs,
\[
 \chi_{\rm o}R_{\rm ext}(\omega)\chi_{\rm s}
\]
has a nonzero rank-one pole of order two at \(\omega_n\).  In the
augmented normalization,
\begin{equation}\label{eq:green-principal}
 \boxed{\;
 \operatorname{Coeff}_{(\omega-\omega_n)^{-2}}
 \bigl(\chi_{\rm o}R_{\rm ext}\chi_{\rm s}\bigr)
 =
 -\frac{\beta_n}{\alpha_n^2}
 (\chi_{\rm o}u_n)\otimes
 (\widetilde u_n\chi_{\rm s}).
 \;}
\end{equation}
\end{theorem}

\begin{proof}
The repeated block inverse is
\[
 \mathbb L^{-1}=
 \begin{pmatrix}
 L^{-1}&-L^{-1}KL^{-1}\\
 0&L^{-1}
 \end{pmatrix}.
\]
At a simple scalar resonance,
\[
 L(\omega)^{-1}
 =\frac{u_n\otimes\widetilde u_n}
 {\alpha_n(\omega-\omega_n)}+O(1).
\]
The upper-right block therefore has principal coefficient
\[
 -\frac{\beta_n}{\alpha_n^2}
 u_n\otimes\widetilde u_n.
\]
The exact transparent conditions identify this finite-interval
coefficient with the cut-off exterior coefficient.  It is nonzero because
\(\beta_n\ne0\), and a nonzero solution of a second-order ODE cannot
vanish on a nonempty open interval.
\end{proof}

The full doubled-state principal coefficient maps the carrier channel to
the Einstein channel and vanishes on the Einstein channel; it is
square-zero on that graded space.  Its projection to a single metric or
radial block is rank one but is not intrinsically a nilpotent operator
after the grading has been forgotten.

\subsection{Mass derivative of the parent response}

At points where the inverses exist,
\[
 \partial_m(\mathcal E+m\mathcal A)^{-1}
 =-(\mathcal E+m\mathcal A)^{-1}
 \mathcal A(\mathcal E+m\mathcal A)^{-1}.
\]
The metric--metric block of the pure-Weyl parent inverse is therefore
\begin{equation}\label{eq:parent-mass-derivative}
 \boxed{\;
 G_{hh}^{\rm W}(\omega)
 =-\frac1{4\alpha_{\rm W}}
 \left.
 \partial_m(\mathcal E+m\mathcal A)^{-1}
 \right|_{m=0}.
 \;}
\end{equation}
No commutativity assumption is used.

If the massive second-order inverse has the local simple-pole form
\[
 R_m(\omega)=
 \frac{P_n(m)}{\omega-\omega_n(m)}+H_n(\omega,m),
\]
then \eqref{eq:parent-mass-derivative} gives
\[
 G_{hh}^{\rm W}(\omega)
 =-\frac1{4\alpha_{\rm W}}
 \left[
 \frac{\omega_n'(0)P_n(0)}
 {(\omega-\omega_n)^2}
 \frac{P_n'(0)}
 {\omega-\omega_n}
 \partial_mH_n(\omega,0)
 \right].
\]
Thus the double pole measures QNM motion; the accompanying simple pole
measures motion of the residue projector.

\begin{corollary}[Isolated resonance contour]
\label{cor:isolated-contour}
For a contour enclosing only the certified local resonance,
\[
 \frac1{2\pi i}\oint
 e^{i\omega t}G_{hh}^{\rm W}(\omega)\,d\omega
 =
 -\frac{e^{i\omega_nt}}{4\alpha_{\rm W}}
 \left[P_n'(0)+it\,\omega_n'(0)P_n(0)\right].
\]
The \(t\)-weighted spatial profile is the ordinary Einstein QNM.  The
constant generalized component contains the non-Einstein mass tangent.
\end{corollary}

This is a local residue theorem, not a global retarded expansion.  The
latter would require a Laplace contour, a closed weighted realization,
control of thresholds and non-pole contributions, and a justified contour
deformation.

\section{Physical source and observation channels}
\label{sec:channels}

\subsection{Metric reconstruction at null infinity}

Let \(E\) be the unit outgoing Einstein line and \(R\) the unit outgoing
carrier line in the exact factor frame.  After factoring the common
retarded phase \(e^{i\omega u}\), their Regge--Wheeler-gauge metric heads
are
\[
 (H_0,H_1)_E=(-r,2r)+O(1),
\]
\[
 (H_0,H_1)_R=
 \left(\frac34r^2-\frac32r,-\frac32r^2\right)+O(1).
\]
Hence
\[
 (H_0,H_1)_R=-\frac34r\,(H_0,H_1)_E+O(E).
\]

In odd radiation gauge,
\[
 h_u\mapsto h_u-i\omega\xi,\qquad
 h_2\mapsto h_2-2\xi,\qquad
 \xi=\frac{h_u}{i\omega},
\]
so \(h_2=2ih_u/\omega\).  Writing \(h_{AB}=h_2X_{AB}\), define the
Bondi-shear observation map by
\[
 \mathcal O_{\mathscr I^+}h
 =\lim_{r\to\infty}\frac{h_{AB}}r.
\]

\begin{theorem}[Null-infinity reconstruction]
\label{thm:null-infinity}
The outgoing Einstein line has nonzero Bondi shear,
\begin{equation}\label{eq:einstein-shear}
 \boxed{\;
 \mathcal O_{\mathscr I^+}E
 =-\frac{2i}{\omega}X_{AB}\ne0.
 \;}
\end{equation}
Its strain has the standard radiative falloff
\[
 \frac{h^E_{AB}}{r^2}
 =-\frac{2i}{\omega r}e^{i\omega u}X_{AB}
 +O(r^{-2}).
\]
The carrier line instead has
\begin{equation}\label{eq:carrier-falloff}
 \boxed{\;
 \frac{h^R_{AB}}{r^2}
 =\frac{3i}{2\omega}e^{i\omega u}X_{AB}
 +O(r^{-1}).
 \;}
\end{equation}
Therefore the time-independent generalized component has weakened
\(O(1)\) strain in this standard odd radiation gauge.  The Einstein
range of the double-pole coefficient is nevertheless observed
nontrivially at future null infinity.
\end{theorem}

\begin{proof}
Substitute the exact outgoing factor-frame series into the algebraic
metric reconstruction.  The odd gauge law then gives
\eqref{eq:einstein-shear} and \eqref{eq:carrier-falloff}.  Adding an
Einstein representative cannot cancel the carrier \(r^2\) coefficient,
which is one power higher.  The trace-free odd tensor harmonic cannot be
removed by a Weyl rescaling; its removal would require a large boundary
diffeomorphism rather than a standard small gauge transformation.
\end{proof}

Combining Theorems~\ref{thm:exterior-green-pole} and
\ref{thm:null-infinity}, the principal Bondi-shear coefficient is
nonzero.  In the parent normalization it is
\[
 \mathcal O_{\mathscr I^+}G_{-2}
 =
 \frac{i\,\omega_n'(0)}
 {2\alpha_{\rm W}\alpha_n\omega_n}
 X_{AB}\otimes\widetilde u_n.
\]
The isolated contour therefore contains an ordinary \(u/r\) radiative
profile multiplied by the adjoint source overlap.

\subsection{A conserved, traceless source with nonzero overlap}

Let \(\ell\ge2\), \(\mu=(\ell-1)(\ell+2)\), and \(\omega\ne0\).  Use the
Martel--Poisson odd source projections and orientation
\(\epsilon_{tr}=1\), \(\epsilon^{tr}=-1\) \cite{martelpoisson2005}.

\begin{theorem}[Conserved source nonannihilation]
\label{thm:source}
For any
\[
 F\in C_c^\infty((2M,\infty)),
\]
define
\[
 P_t=0,\qquad
 P_r=\frac{\mu F}{2i\omega rf},\qquad
 P=\frac{\partial_r(rF)}{2i\omega}.
\]
Then
\[
 \nabla_aP^a+\frac2r r_aP^a-\frac{\mu}{r^2}P=0,
\]
and the Cunningham--Price--Moncrief source is
\[
 S_{\rm odd}
 =-\frac{2r}{\mu}\epsilon^{ab}\nabla_aP_b
 =\frac{F}{f}.
\]
The associated odd stress tensor
\[
 T^{aB}=\frac{P^a}{16\pi r^2}X^B,\qquad
 T^{AB}=\frac{P}{8\pi r^4}X^{AB},
\]
with all other harmonic components zero, is smooth and radially compact,
and satisfies
\[
 \boxed{\;
 \nabla_\mu T^{\mu\nu}=0,\qquad
 T^\mu{}_\mu=0.
 \;}
\]
The complexified conserved-traceless odd source map is therefore onto the
smooth compact reduced radial forcings.

Moreover, choose a nonnegative nonzero
\(\eta\in C_c^\infty((2M,\infty))\) supported where the adjoint state
\(\widetilde u_n\) is nonzero and set
\[
 F=\eta\,\overline{\widetilde u_n}.
\]
Then the augmented pairing has no endpoint term and
\[
 \boxed{\;
 \langle\widetilde u_n,F\rangle_{\rm aug}
 =\int\eta|\widetilde u_n|^2\,dr_*>0.
 \;}
\]
Thus a smooth compact conserved and traceless odd source excites the
certified double pole.
\end{theorem}

\begin{proof}
Since \(\sqrt{-\det g_{ab}}=1\),
\[
 P^r=fP_r=\frac{\mu F}{2i\omega r},
\qquad
 \nabla_aP^a=\partial_rP^r.
\]
Direct substitution proves the reduced conservation identity and
\(fS_{\rm odd}=F\).  The remaining spacetime conservation equations
vanish by the divergence-free odd vector harmonic, and tracelessness
follows from \(T^{ab}=0\) and
\(\Omega_{AB}X^{AB}=0\).  Compact support removes all endpoint terms from
the adjoint pairing, and the final integral is strictly positive.
\end{proof}

\begin{remark}[What the source theorem does not say]
The theorem constructs a conformally admissible external source and proves
nonannihilation.  It does not identify a specified material trajectory or
an energy-condition-satisfying matter model.  In particular, a standard
massive point-particle stress tensor is not trace-free and cannot be
coupled directly to the Bach equation without a conformal completion or a
compensator.
\end{remark}

\section{Scope, evidence, and open problem}
\label{sec:scope}

\subsection{The remaining analytic problem}

The exterior cut-off theorem is already a differential-response result on
the Schwarzschild exterior.  The remaining causal problem is narrower
but substantive.  A global realization must provide:
\begin{enumerate}
\item closed weighted domains implementing horizon-ingoing and
  differentiated outgoing Jost conditions;
\item a Fredholm pencil of index zero with a fixed analytic domain;
\item bounded reconstruction through the physical gauge quotient;
\item a causal Laplace contour and a justified deformation through the
  QNM disk;
\item control of threshold, branch-cut, and non-pole contributions.
\end{enumerate}

The weakened falloff in \eqref{eq:carrier-falloff} explains why this is
not routine.  A standard outgoing asymptotically flat function space may
contain the Einstein root but exclude its mass tangent.  One may instead
use an enlarged tangent domain or keep the bulk space fixed and encode the
differentiated endpoint germ as a finite-dimensional boundary component.
This paper uses the latter construction locally but does not complete its
global causal implementation.

\subsection{Certificate dependency}

\begin{center}
\footnotesize
\begin{tabular}{p{0.29\linewidth}p{0.44\linewidth}p{0.17\linewidth}}
\toprule
Certificate family & Imported statement & Type\\
\midrule
Axial triangular preflight
& Filtration, factor gauge, and reconstruction rank
& Exact\\
Projective cocycle
& Cocycle, primitive, and rational nonsplitting obstruction
& Exact\\
Critical mass jet
& Scalar mass direction and parent inverse identity
& Exact\\
Analytic continuation
& Analytic horizon and infinity Jost germs on the QNM disk
& Exact gate\\
Full Evans contour
& Zero-free contour and winding \(+1\)
& Validated\\
Local selector
& Nonzero projective tangent on the root disk
& Validated\\
Spin-one local unit
& Nonzero spin-one mismatch
& Validated\\
Fredholm promotion
& Finite-interval Poisson/trace reduction
& Exact reduction\\
Null-infinity reconstruction
& Metric heads, radiation gauge, and Bondi shear
& Exact\\
Conserved source
& Conservation, tracelessness, and adjoint overlap
& Exact\\
\bottomrule
\end{tabular}
\end{center}

The generated claim map pins the authority files and their hashes.  Its
verifier independently checks the displayed cocycle identity, the
mass-parameter normalization, the Smith and Green formulas, the
reconstruction falloffs, the source conservation law, and explicit
non-results.  Missing data, a stale hash, an unsupported domain, or a
failed mutation is a rejection rather than a partial pass.

\paragraph{Explicit non-results.}
This paper does not establish a global causal spacetime resolvent, a
complete QNM expansion, a late-time asymptotic formula, a retarded
inverse-Laplace deformation, a time-domain instability, finite Bondi flux
for the generalized constant component, a standard asymptotically flat
domain containing that component, excitation by a specified
astrophysical source, detector sensitivity, an overtone tower, an
all-multipole Bach nonsplitting theorem, or a quantum positivity/no-go
theorem.  The polynomially weighted term is proved only for the isolated
local resonance contour.

\section{Conclusion}

The repeated Regge--Wheeler factor in axial Schwarzschild pure Weyl
gravity is not a direct sum.  It is a non-split self-extension with the
elementary projective class
\[
 [\mathcal I_{\rm Bach}]
 =\frac{i\omega}{2}[1-2/r].
\]
That direction is exactly the critical transverse-traceless mass jet.
The linearly growing rational equivalence gauge is compatible with the
differentiated physical Jost germs because the massive Coulomb exponent
has vanishing first derivative.

At the certified damped Schwarzschild mode, the scalar QNM is simple but
the extension selector is nonzero.  The full connection is therefore
defective.  Its geometric root is Einstein; its generalized root has a
nonzero Ricci-carrier quotient and is the massive-QNM tangent modulo the
ordinary root.  The same invariant controls the nonzero rank-one
second-order pole of the exterior cut-off outgoing Green operator.

The leading pole survives the two physically relevant channel tests.
Its Einstein range has nonzero Bondi shear, and a smooth compact
conserved, traceless odd source can have strictly nonzero adjoint overlap.
The enhanced isolated-contour profile is therefore an ordinary outgoing
Einstein QNM multiplied by time, even though the constant generalized
component has weakened asymptotic falloff.

This gives a Schwarzschild counterpart, in asymptotically flat spacetime,
of the critical-gravity logarithmic-partner mechanism.  The open step is not the
existence of the defective radial response; it is the construction of a
global causal domain that accommodates the differentiated endpoint state
and supports a justified retarded contour deformation.

\begin{thebibliography}{9}

\bibitem{lupope2011}
H.~L\"u and C.~N.~Pope,
\newblock Critical gravity in four dimensions,
\newblock \emph{Physical Review Letters} \textbf{106} (2011), 181302,
\newblock \href{https://doi.org/10.1103/PhysRevLett.106.181302}
{doi:10.1103/PhysRevLett.106.181302}.

\bibitem{bergshoeff2011}
E.~A.~Bergshoeff, O.~Hohm, J.~Rosseel, and P.~K.~Townsend,
\newblock Modes of log gravity,
\newblock \emph{Physical Review D} \textbf{83} (2011), 104038,
\newblock \href{https://doi.org/10.1103/PhysRevD.83.104038}
{doi:10.1103/PhysRevD.83.104038}.

\bibitem{martelpoisson2005}
K.~Martel and E.~Poisson,
\newblock Gravitational perturbations of the Schwarzschild spacetime: A
practical covariant and gauge-invariant formalism,
\newblock \emph{Physical Review D} \textbf{71} (2005), 104003,
\newblock \href{https://doi.org/10.1103/PhysRevD.71.104003}
{doi:10.1103/PhysRevD.71.104003}.

\end{thebibliography}

\end{document}
